% GENERATED FILE. Edit canonical lessons, then run npm run generate:books.
% canonical-source-hash: fnv1a64:13b5055063fe3f24
% canonical-lessons: TA-C409-kuli, TA-C409-pallam, TA-C409-katavai, TA-C409-kurukkuvali, TA-C409-cikku

\chapter{Potholes and Shortcuts}
\label{ch:ta-potholes-and-shortcuts}
\hlchaptermodality{409}

\begin{chapteropening}
\textbf{By the end of this chapter:} \emph{I can talk about pits, potholes, a level crossing and a shortcut, and getting stuck.}

Complete the last lesson of chapter 409: I can talk about pits, potholes, a level crossing and a shortcut, and getting stuck.
\end{chapteropening}

\section[\texorpdfstring{\ta{குழி} (kuḻi)}{குழி (kuḻi)}]{\ta{குழி} (kuḻi) — a pit; a hole in the ground}

\label{lesson:TA-C409-kuli}



\begin{quote}
Before the new one: say the Tamil for to overturn, then the Tamil for to spin.
\end{quote}

\subsection*{You'll want to know: \ta{குழி}}
\textbf{\ta{குழி}} — \emph{kuḻi} — \textquotedblleft{}a pit; a hole in the ground\textquotedblright{}.

\begin{quote}
\ta{சாலையில்} \ta{ஒரு} \ta{பெரிய} \ta{குழி} \ta{இருந்தது}.
\end{quote}

\begin{quote}
\emph{cālaiyil oru periya kuḻi iruntatu.}
\end{quote}

\begin{quote}
There was a big pit in the road.
\end{quote}

The plural is \ta{குழிகள்}.

\begin{cousinweb}
\textbf{\ta{குழி}} (\emph{kuḻi}) is native Dravidian (DEDR 1818), \textquotedblleft{}a pit, a hole, a hollow\textquotedblright{}.
\end{cousinweb}

\begin{grammarlens}[title={What you've built}]
One more word for getting around.
\end{grammarlens}

\subsection*{Your turn}
\begin{itemize}
\raggedright
  \item \emph{Say it:} \emph{kuḻi}
  \item \emph{Say it:} \emph{kuḻi}, once more
  \item \emph{Say it:} \emph{cālaiyil oru periya kuḻi iruntatu.} — There was a big pit in the road.
  \item \emph{Recall:} say the Tamil for to overturn, then the Tamil for to spin, then say \emph{kuḻi} again
\end{itemize}

\begin{tcolorbox}[breakable,colback=teal!4,colframe=teal!35!black,title={Before you move on}]
What does \ta{குழி} mean? (\textquotedblleft{}A pit; a hole in the ground\textquotedblright{}.) Say it once more.
\end{tcolorbox}
\section[\texorpdfstring{\ta{பள்ளம்} (paḷḷam)}{பள்ளம் (paḷḷam)}]{\ta{பள்ளம்} (paḷḷam) — a pothole; a hollow}

\label{lesson:TA-C409-pallam}



\begin{quote}
Before the new one: say the Tamil for to spin, then the Tamil for a pit.
\end{quote}

\subsection*{You'll want to know: \ta{பள்ளம்}}
\textbf{\ta{பள்ளம்}} — \emph{paḷḷam} — \textquotedblleft{}a pothole; a hollow\textquotedblright{}.

\begin{quote}
\ta{மழைக்குப்} \ta{பிறகு} \ta{பள்ளங்களில்} \ta{தண்ணீர்} \ta{இருந்தது}.
\end{quote}

\begin{quote}
\emph{maḻaikkup piṟaku paḷḷaṅkaḷil taṇṇīr iruntatu.}
\end{quote}

\begin{quote}
After the rain there was water in the potholes.
\end{quote}

Plural: \ta{பள்ளங்கள்}, \emph{paḷḷaṅkaḷ}.

A -\ta{ம்} noun: plural \ta{பள்ளங்கள்}. \ta{பள்ளத்தாக்கு}, a valley, is built on it.

\begin{cousinweb}
\textbf{\ta{பள்ளம்}} (\emph{paḷḷam}) is native Dravidian (DEDR 4016), \textquotedblleft{}lowness, low land, a hollow, a ditch\textquotedblright{}.
\end{cousinweb}

\begin{grammarlens}[title={What you've built}]
One more word for getting around.
\end{grammarlens}

\subsection*{Your turn}
\begin{itemize}
\raggedright
  \item \emph{Say it:} \emph{paḷḷam}
  \item \emph{Say it:} \emph{paḷḷam}, once more
  \item \emph{Say it:} \emph{maḻaikkup piṟaku paḷḷaṅkaḷil taṇṇīr iruntatu.} — After the rain there was water in the potholes.
  \item \emph{Recall:} say the Tamil for to spin, then the Tamil for a pit, then say \emph{paḷḷam} again
\end{itemize}

\begin{tcolorbox}[breakable,colback=teal!4,colframe=teal!35!black,title={Before you move on}]
What does \ta{பள்ளம்} mean? (\textquotedblleft{}A pothole; a hollow\textquotedblright{}.) Say it once more.
\end{tcolorbox}
\section[\texorpdfstring{\ta{கடவை} (kaṭavai)}{கடவை (kaṭavai)}]{\ta{கடவை} (kaṭavai) — a level crossing; a crossing place}

\label{lesson:TA-C409-katavai}



\begin{quote}
Before the new one: say the Tamil for a pit, then the Tamil for a pothole.
\end{quote}

\subsection*{You'll want to know: \ta{கடவை}}
\textbf{\ta{கடவை}} — \emph{kaṭavai} — \textquotedblleft{}a level crossing; a crossing place\textquotedblright{}.

\begin{quote}
\ta{ரயில்} \ta{போகும்வரை} \ta{நாங்கள்} \ta{கடவையில்} \ta{காத்திருந்தோம்}.
\end{quote}

\begin{quote}
\emph{rayil pōkumvarai nāṅkaḷ kaṭavaiyil kāttiruntōm.}
\end{quote}

\begin{quote}
We waited at the level crossing until the train went by.
\end{quote}

From \ta{கட}, to cross, + -\ta{வை}: a place for crossing.

\begin{cousinweb}
\textbf{\ta{கடவை}} (\emph{kaṭavai}) is built from \ta{கட}, \textquotedblleft{}to cross, to pass through\textquotedblright{} (DEDR 1109), with a noun ending: a place to cross.
\end{cousinweb}

\begin{grammarlens}[title={What you've built}]
One more word for getting around.
\end{grammarlens}

\subsection*{Your turn}
\begin{itemize}
\raggedright
  \item \emph{Say it:} \emph{kaṭavai}
  \item \emph{Say it:} \emph{kaṭavai}, once more
  \item \emph{Say it:} \emph{rayil pōkumvarai nāṅkaḷ kaṭavaiyil kāttiruntōm.} — We waited at the level crossing until the train went by.
  \item \emph{Recall:} say the Tamil for a pit, then the Tamil for a pothole, then say \emph{kaṭavai} again
\end{itemize}

\begin{tcolorbox}[breakable,colback=teal!4,colframe=teal!35!black,title={Before you move on}]
What does \ta{கடவை} mean? (\textquotedblleft{}A level crossing; a crossing place\textquotedblright{}.) Say it once more.
\end{tcolorbox}
\section[\texorpdfstring{\ta{குறுக்குவழி} (kuṟukkuvaḻi)}{குறுக்குவழி (kuṟukkuvaḻi)}]{\ta{குறுக்குவழி} (kuṟukkuvaḻi) — a shortcut}

\label{lesson:TA-C409-kurukkuvali}



\begin{quote}
Before the new one: say the Tamil for a pothole, then the Tamil for a level crossing.
\end{quote}

\subsection*{You'll want to know: \ta{குறுக்குவழி}}
\textbf{\ta{குறுக்குவழி}} — \emph{kuṟukkuvaḻi} — \textquotedblleft{}a shortcut\textquotedblright{}.

\begin{quote}
\ta{குறுக்குவழியில்} \ta{போனதால்} \ta{சீக்கிரம்} \ta{வந்தோம்}.
\end{quote}

\begin{quote}
\emph{kuṟukkuvaḻiyil pōṉatāl cīkkiram vantōm.}
\end{quote}

\begin{quote}
We arrived early because we took a shortcut.
\end{quote}

\ta{குறுக்கு}, across, + \ta{வழி}, way.

\begin{cousinweb}
\textbf{\ta{குறுக்குவழி}} (\emph{kuṟukkuvaḻi}) is built from \ta{குறுக்கு}, \textquotedblleft{}to shorten; a short distance, across\textquotedblright{} (DEDR 1851), and \ta{வழி}, \textquotedblleft{}a way\textquotedblright{} (DEDR 5297).
\end{cousinweb}

\begin{grammarlens}[title={What you've built}]
One more word for getting around.
\end{grammarlens}

\subsection*{Your turn}
\begin{itemize}
\raggedright
  \item \emph{Say it:} \emph{kuṟukkuvaḻi}
  \item \emph{Say it:} \emph{kuṟukkuvaḻi}, once more
  \item \emph{Say it:} \emph{kuṟukkuvaḻiyil pōṉatāl cīkkiram vantōm.} — We arrived early because we took a shortcut.
  \item \emph{Recall:} say the Tamil for a pothole, then the Tamil for a level crossing, then say \emph{kuṟukkuvaḻi} again
\end{itemize}

\begin{tcolorbox}[breakable,colback=teal!4,colframe=teal!35!black,title={Before you move on}]
What does \ta{குறுக்குவழி} mean? (\textquotedblleft{}A shortcut\textquotedblright{}.) Say it once more.
\end{tcolorbox}
\section[\texorpdfstring{\ta{சிக்கு} (cikku)}{சிக்கு (cikku)}]{\ta{சிக்கு} (cikku) — to get stuck; to be caught}

\label{lesson:TA-C409-cikku}



\begin{quote}
Before the new one: say the Tamil for a level crossing, then the Tamil for a shortcut.
\end{quote}

\subsection*{You'll want to know: \ta{சிக்கு}}
\textbf{\ta{சிக்கு}} — \emph{cikku} — \textquotedblleft{}to get stuck; to be caught\textquotedblright{}.

\begin{quote}
\ta{எங்கள்} \ta{கார்} \ta{சேற்றில்} \ta{சிக்கியது}.
\end{quote}

\begin{quote}
\emph{eṅkaḷ kār cēṟṟil cikkiyatu.}
\end{quote}

\begin{quote}
Our car got stuck in the mud.
\end{quote}

Past (I …): \ta{சிக்கினேன்}, \emph{cikkiṉēṉ}.

Where one is stuck takes -\ta{இல்}: \ta{சேற்றில்} \ta{சிக்கு}; often with \ta{கொள்}: \ta{சிக்கிக்கொண்டேன்}, I got caught.

\begin{cousinweb}
\textbf{\ta{சிக்கு}} (\emph{cikku}) is native Dravidian (DEDR 2498), \textquotedblleft{}to become entangled, to be caught\textquotedblright{}.
\end{cousinweb}

\begin{grammarlens}[title={What you've built}]
One more word for getting around.
\end{grammarlens}

\subsection*{Your turn}
\begin{itemize}
\raggedright
  \item \emph{Say it:} \emph{cikku}
  \item \emph{Say it:} \emph{cikku}, once more
  \item \emph{Say it:} \emph{eṅkaḷ kār cēṟṟil cikkiyatu.} — Our car got stuck in the mud.
  \item \emph{Recall:} say the Tamil for a level crossing, then the Tamil for a shortcut, then say \emph{cikku} again
\end{itemize}

\begin{tcolorbox}[breakable,colback=teal!4,colframe=teal!35!black,title={Before you move on}]
What does \ta{சிக்கு} mean? (\textquotedblleft{}To get stuck; to be caught\textquotedblright{}.) Say it once more.
\end{tcolorbox}
